\section{Introdução}
\label{sec:intro}

Ataques distribuídos de negação de serviço (DDoS) continuam entre as ameaças mais frequentes à
disponibilidade de serviços em rede, e os relatórios de operadoras de mitigação registram novos
vetores a cada trimestre~\cite{cloudflare2026}. Detectores supervisionados dependem de tráfego
rotulado, que envelhece rápido e não cobre vetores novos. Por isso há interesse em detectores
\emph{online} e \emph{não supervisionados}, que aprendem o perfil do tráfego legítimo durante a
operação e apontam o que foge dele~\cite{mirsky2018kitsune}. Métodos de isolamento, como o
Isolation Forest~\cite{liu2008iforest} e sua variante para fluxos, as Half-Space
Trees~\cite{tan2011hst}, são a base usual desses detectores, pois têm custo constante por instância.

O \modus~\cite{modus2026}, publicado no SBSeg 2026, segue essa linha. O método combina um detector
de mudança de conceito (ADWIN), um HST e um limiar baseado em mediana e MAD, e reporta F1 de
99,62\%, 83,79\% e 81,24\% em três bases públicas. O código dos autores não está disponível. Este
artigo relata uma replicação independente, feita em duas etapas na disciplina de Auditoria e
Segurança de Sistemas de Informação do CIn/UFPE.

Na Etapa~1 reimplementamos o método a partir do texto e comparamos os resultados com os do artigo.
Dois dos três números se reproduziram, e a análise apontou fragilidades: a \emph{flag} de mudança
de conceito não altera as predições, o HST sem normalização separa por protocolo e não por volume,
e o limiar depende de um piso ajustado por base. A etapa terminou com quatro melhorias propostas e
três defesas indicadas como trabalho futuro.

Na Etapa~2, objeto principal deste artigo, implementamos as melhorias. Para controlar a janela de
referência do HST, reescrevemos o algoritmo em \texttt{numpy}, bit a bit igual ao do MOA. A
reescrita revelou um defeito que muda a leitura da Etapa~1: o \texttt{HSTrees} do MOA inclui o
atributo de classe entre as dimensões de corte. Quando o fluxo carrega o rótulo verdadeiro, o
detector ``não supervisionado'' usa o rótulo. A réplica dos 99,62\% no ASEADOS-SDN-IoT só existe
com esse vazamento.

Removido o vazamento, o detector falha por outro motivo, mais geral. O treino seletivo, que só
aprende com o que foi predito como benigno, não protege o modelo: basta o início do ataque passar
pelo limiar para que o \emph{flood}, por ser denso, domine a massa do HST e vire o perfil de
referência. A partir daí o ataque é o tráfego \emph{mais normal} do fluxo.

\subsection*{Contribuições}
\begin{itemize}
\item \textbf{Replicação independente} do \modus\ nas três bases do artigo, com a reconstrução
exata do número de instâncias de cada uma e a identificação de inconsistências nos rótulos do
CTU-13 e na composição do Edge-IIoTset (Seção~\ref{sec:replication}).
\item \textbf{Vazamento de rótulo no \texttt{HSTrees} do MOA}: descrição do mecanismo, teste direto
e medição do efeito por base. No ASEADOS o F1 vai de 99,73\% para 0,00\% (Seção~\ref{sec:leak}).
\item \textbf{HST em \texttt{numpy} equivalente ao do MOA}, com diferença máxima de
$2{,}2\times10^{-16}$ nos scores, sem a dimensão do rótulo e com controle da janela de referência
(Seção~\ref{sec:hstnp}).
\item \textbf{Medição do envenenamento por \emph{flood}}: a fração do ataque usada no treino do
modelo, reportada em todas as rodadas (Seção~\ref{sec:results}).
\item \textbf{\modusp}: quatro defesas causais e sem rótulos, avaliadas com uma única configuração
fixada a priori, com ablação, sementes e sensibilidade (Seções~\ref{sec:method}
e~\ref{sec:results}).
\end{itemize}

Os resultados são mistos, e os reportamos assim. O \modusp\ leva o F1 de 0,08\% para 62,99\% no
ASEADOS e de 76,66\% para 91,17\% no Edge-IIoTset com 14 tipos de ataque. No CTU-13 o
envenenamento praticamente desaparece, mas o F1 a priori fica em 0,01\%, porque o score do
\emph{flood} fica milésimos abaixo do limiar. Todos os números deste artigo vêm de rodadas salvas
no repositório do trabalho, com configuração, \emph{log} e métricas.

O restante do texto está organizado assim. A Seção~\ref{sec:background} apresenta os fundamentos e
o \modus. A Seção~\ref{sec:replication} resume a replicação e descreve o vazamento. A
Seção~\ref{sec:method} detalha o \modusp, e a Seção~\ref{sec:setup}, a metodologia. A
Seção~\ref{sec:results} traz os resultados, a Seção~\ref{sec:discussion} os discute, e a
Seção~\ref{sec:conclusion} conclui.
